\chapter{Локальность, непрерывные состояния и импульс}

\lectureoverview{Запутанность может создаваться и устраняться обратимой динамикой, однако измерения становятся фактически необратимыми, когда часть запутанной системы уходит из-под нашего контроля. Это различие приводит к точной проверке локальности: Боб может применять в своей окрестности любую унитарную операцию, но редуцированная матрица плотности Алисы, а значит и любая статистика, доступная ей одной, остается неизменной. Затем мысленный эксперимент с классическим моделированием вычленяет то, что действительно неклассично. Один спин легко имитировать, и центральный компьютер может имитировать запутанную пару, но разнесенные локальные машины не могут воспроизвести все запутанные корреляции без недозволенной связи или белл-локальной модели со скрытыми переменными, которая в конечном счете оказывается несостоятельной. Затем лекция переходит от спинов к частице на прямой. Дискретные суммы превращаются в интегралы, дельты Кронекера — в дельты Дирака, а волновые функции — в векторы в бесконечномерном пространстве. Умножение на \(x\) дает оператор координаты; оператор дифференцирования антиэрмитов; умножение его на \(-i\) дает эрмитов оператор импульса. Заключительное вычисление, \([x,p]=i\hbar\), вновь связывает квантовые коммутаторы с классическими скобками Пуассона и указывает на волновые пакеты, преобразования Фурье и Schr\"odinger-эволюцию.}

\section{Запутанность — не необратимость}

Начнем с различия, которое легко упустить в обсуждениях измерения. Предположим, что два спина начинают в состоянии-произведении, взаимодействуют, затем расходятся и оказываются в запутанном состоянии. Должна ли запутанность сохраняться навсегда? Нет. Взаимодействие было унитарным, поэтому другое подходящее взаимодействие может его обратить. Можно отразить спины, снова свести их вместе и организовать обратную эволюцию, возвращающую пару в состояние-произведение. Не существует закона, согласно которому сама запутанность может только возрастать.

С операциональной точки зрения ситуация меняется, когда один из членов пары отбрасывают. Если электрон улетает в неконтролируемую среду, то остающийся поблизости партнер по-прежнему с чем-то запутан, но недостающие степени свободы уже недоступны для обратной операции. В принципе далекий электрон можно было бы отразить и вернуть; на практике он потерян. Именно в этом смысле измерение необратимо \emph{для нас}: корреляции распространились на степени свободы, которыми мы не управляем.

\begin{classroomqa}
\audiencequestion{Если две системы стали запутанными, могут ли они позже утратить запутанность?}
\lecturerresponse{Да. Унитарная эволюция обратима. Второе взаимодействие может реализовать преобразование, обратное первому, и восстановить состояние-произведение. Необратимость появляется, когда релевантные степени свободы отбрасываются или поглощаются неконтролируемой средой, а не из одной лишь запутанности.}
\end{classroomqa}

\begin{classroomqa}
\audiencequestion{Тогда почему завершившееся измерение кажется необратимым?}
\lecturerresponse{Потому что измеряемая система оказалась скоррелированной с прибором, наблюдателем и, в конечном счете, с большой средой. Чтобы обратить полное состояние, потребовался бы контроль над всеми этими записями. Микроскопическая эволюция может оставаться унитарной, тогда как операция недоступна в редуцированном описании наблюдателя.}
\end{classroomqa}

\section{Состояние, доступное Алисе}

Разобьем составную систему на степени свободы Алисы \(a\) и степени свободы Боба \(b\), а коэффициенты чистого состояния запишем как
\begin{equation}
\Psi(a,b)=\braket{a,b}{\Psi}.
\end{equation}
У Алисы нет доступа к лаборатории Боба. Поэтому ее полное локальное описание, вообще говоря, не является волновой функцией. Это редуцированная матрица плотности
\begin{equation}
\rho_A(a',a)
=
\sum_b \Psi^*(a',b)\Psi(a,b).
\label{eq:l8-reduced-density}
\end{equation}
Суммирование устраняет базисную метку Боба. То, что остается, имеет один бра-индекс и один кет-индекс Алисы и содержит каждое распределение вероятностей, которое Алиса может получить в экспериментах, полностью проводимых в ее окрестности.

\begin{figure}[htbp]
\centering
\includegraphics[width=0.88\textwidth]{lecture_08_figure_01.png}
\caption{Редуцированная матрица плотности Алисы, полученная суммированием совместной волновой функции по базисной метке Боба.}
\lectureframe{00:04:34}
\end{figure}

Здесь «окрестность» Алисы и «окрестность» Боба следует понимать широко. Каждая из них может включать приборы, магнитные поля, техников и любые другие степени свободы, с которыми данный наблюдатель может взаимодействовать. Утверждение о локальности не состоит в том, что Боб не может изменить свою собственную систему. Он может делать это многими способами. Утверждение состоит в том, что, пока он остается удаленным, ничто из того, что он делает локально, не изменяет формулу~\eqref{eq:l8-reduced-density}.

\section{Локальная унитарная операция Боба}

Сначала отвлечемся от Алисы и опишем только действие Боба. Если Боб изменяет магнитное поле или иным образом меняет свой локальный гамильтониан, коэффициенты его волновой функции претерпевают унитарное преобразование,
\begin{equation}
\psi_M(b)
=
\sum_{b'}U_{bb'}\psi(b').
\label{eq:l8-bob-unitary}
\end{equation}
В описании составной системы индекс Алисы при этом не затрагивается:
\begin{equation}
\Psi_M(a,b)
=
\sum_{b'}U_{bb'}\Psi(a,b').
\label{eq:l8-composite-unitary}
\end{equation}
Комплексное сопряжение переводит преобразование кет-вектора в соответствующее преобразование бра-вектора,
\begin{equation}
\Psi_M^*(a',b)
=
\sum_{b''}\Psi^*(a',b'')U^\dagger_{b''b}.
\label{eq:l8-composite-bra}
\end{equation}

\begin{figure}[htbp]
\centering
\includegraphics[width=0.88\textwidth]{lecture_08_figure_02.png}
\caption{Унитарная матрица действует на кет-индекс Боба, а эрмитово-сопряженная матрица — на соответствующий бра-индекс.}
\lectureframe{00:12:54}
\end{figure}

Теперь пересчитаем матрицу плотности Алисы:
\begin{align}
\rho_{A,M}(a',a)
&=
\sum_b\Psi_M^*(a',b)\Psi_M(a,b) \notag\\
&=
\sum_{b,b',b''}
\Psi^*(a',b'')
U^\dagger_{b''b}U_{bb'}
\Psi(a,b') \notag\\
&=
\sum_{b',b''}
\Psi^*(a',b'')
(U^\dagger U)_{b''b'}
\Psi(a,b') \notag\\
&=
\sum_{b',b''}
\Psi^*(a',b'')\delta_{b''b'}\Psi(a,b') \notag\\
&=
\sum_{b'}\Psi^*(a',b')\Psi(a,b') \notag\\
&=
\rho_A(a',a).
\label{eq:l8-locality-proof}
\end{align}

\begin{figure}[htbp]
\centering
\includegraphics[width=0.88\textwidth]{lecture_08_figure_03.png}
\caption{Вычисление локальности на доске: \(U^\dagger U\) сводится к единичному оператору в выражении для редуцированной матрицы плотности Алисы.}
\lectureframe{00:14:27}
\end{figure}

Все доказательство сводится к свертке
\begin{equation}
\sum_bU^\dagger_{b''b}U_{bb'}
=(U^\dagger U)_{b''b'}
=\delta_{b''b'}.
\label{eq:l8-unitarity-contraction}
\end{equation}
Никакого предположения о факторизации не использовалось. Формула~\eqref{eq:l8-locality-proof} верна как тогда, когда совместное состояние является состоянием-произведением, так и тогда, когда оно максимально запутано. Следовательно, любое среднее значение, относящееся только к Алисе,
\begin{equation}
\expect{L_A}
=
\operatorname{Tr}_A(\rho_A L_A),
\end{equation}
инвариантно относительно локальной унитарной операции Боба.

\begin{classroomqa}
\audiencequestion{Если Боб что-то делает со своим спином, почему это уже не является нарушением унитарности?}
\lecturerresponse{Делать что-то не означает применять неунитарное отображение. Боб может повернуть магнит, изменить локальный гамильтониан или позволить своему спину прецессировать. Все это обычные унитарные эволюции. Они изменяют его собственное состояние и вероятности его собственных исходов, но приведенное выше вычисление показывает, что они не изменяют редуцированное состояние Алисы.}
\end{classroomqa}

\begin{classroomqa}
\audiencequestion{Значит ли это, что квантовая механика удовлетворяет локальности?}
\lecturerresponse{Она удовлетворяет той операциональной локальности, которая используется здесь: Боб не может, действуя только в своей окрестности, повлиять на статистику измерений, доступных Алисе. Запутанные совместные корреляции остаются неклассическими, но не дают управляемого сигнала.}
\end{classroomqa}

\subsection{Почему важна унитарность}

Унитарность здесь не была просто украшением. Именно она дает единичный оператор в формуле~\eqref{eq:l8-unitarity-contraction}. В прежних терминах унитарность выражает сохранение различий: два различимых вектора состояния остаются различимыми при обратимой эволюции. Если бы в распоряжении Боба была хотя бы слегка неунитарная детерминированная эволюция, \(U^\dagger U\) уже не было бы единичным оператором. Тогда редуцированная матрица плотности Алисы могла бы немедленно измениться, давая Бобу средство передавать сигналы быстрее света.

Эта тесная связь — одна из причин, по которым квантовую механику трудно изменять по частям. Измените закон эволюции — и вместе с ним могут нарушиться, казалось бы, далекие принципы. Локальность, сохранение вероятности и обратимая эволюция состояния составляют части единой структуры.

\begin{classroomqa}
\audiencequestion{Как неунитарная операция позволила бы передавать сигнал быстрее света?}
\lecturerresponse{Боб выбирал бы, применять ее или нет. Если бы отображение изменяло \(\rho_A\), Алиса могла бы различить эти два выбора, повторяя локальные измерения на ансамбле. Изменение происходило бы без распространения сигнала от Боба, и именно это предотвращает унитарное сокращение.}
\end{classroomqa}

\section{Коллапс меняет условные состояния, а не удаленные маргинальные распределения}

Пример с делителем пучка делает это беспокойство особенно наглядным. Одиночный фотон подготавливают в суперпозиции двух путей: один ведет к Алисе, другой — к Бобу. Если детектор Алисы срабатывает, условное состояние приписывает нулевую амплитуду более позднему срабатыванию на пути Боба. Возникает искушение сказать, что Алиса мгновенно изменила фотон Боба.

Однако то, что Боб может наблюдать без зарегистрированного у Алисы результата, иное. До получения какого-либо сообщения он должен усреднить по возможным исходам Алисы. Это безусловное среднее описывается его редуцированной матрицей плотности и не изменяется. Срабатывание детектора Алисы меняет то, какой условный подансамбль описывает данный прогон после того, как ее результат становится известным; оно не меняет локального маргинального распределения Боба. Различие между
\begin{equation}
\rho_B
\qquad\text{and}\qquad
\rho_{B\mid A=\text{click}}
\end{equation}
и есть различие между отсутствием сигнализации и полезными корреляциями.

\begin{classroomqa}
\audiencequestion{Когда один детектор срабатывает, разве волновая функция на другом плече не коллапсирует мгновенно в нуль?}
\lecturerresponse{Состояние, обусловленное этим срабатыванием, изменяется, но удаленный наблюдатель не располагает этой информацией, пока не придет обычное сообщение. Безусловная матрица плотности Боба и вся его усредненная локальная статистика остаются неизменными. Коллапс нельзя использовать как телеграф.}
\end{classroomqa}

\begin{classroomqa}
\audiencequestion{Не является ли эта путаница отчасти смешением состояния и информации о состоянии?}
\lecturerresponse{Такая формулировка наводит на мысль, но она всё ещё слишком расплывчата. Конкретное различие проходит между совместным состоянием, редуцированным состоянием и состоянием, обусловленным сообщённым результатом. Надёжнее держать эти три математических объекта раздельно, чем пытаться решить вопрос одним лишь словом информация.}
\end{classroomqa}

Тем не менее, беспокойство Эйнштейна--Подольского--Розена сохраняется. Синглет может давать нам максимальное знание о \emph{целом}, тогда как ни одна из составляющих не имеет чистого локального состояния. Этот переворот классического ожидания---знание целого без приписывания частям определённых чистых состояний---не является эффектом передачи сигнала, но он действительно необычен. Чтобы резче выявить эту особенность, превратим квантовую механику в задачу программирования.

\section{Классический симулятор одного спина}

Нарисуйте на экране компьютера спиновый прибор. Пользователь может задавать его ориентацию, нажимать кнопку измерения \(M\) и считывать \(+1\) или \(-1\). Внутри классического компьютера хранятся две комплексные амплитуды,
\begin{equation}
\alpha_u,\qquad \alpha_d,
\qquad
|\alpha_u|^2+|\alpha_d|^2=1.
\label{eq:l8-one-spin-amplitudes}
\end{equation}
Между измерениями компьютер интегрирует соответствующее уравнение Schr\"odinger. Смоделированное магнитное поле может заставить амплитуды прецессировать точно так же, как прецессировало бы состояние реального спина.

\begin{figure}[htbp]
\centering
\includegraphics[width=0.88\textwidth]{lecture_08_figure_04.png}
\caption{Симулятор одного спина: ориентируемый прибор и кнопка измерения, соединённые с компьютером, хранящим две амплитуды.}
\lectureframe{00:27:47}
\end{figure}

Когда пользователь нажимает \(M\), генератор псевдослучайных чисел выбирает результат с вероятностями Борна
\begin{equation}
P(+1)=|\alpha_u|^2,
\qquad
P(-1)=|\alpha_d|^2.
\label{eq:l8-one-spin-born}
\end{equation}
Компьютер выводит выбранный знак, а затем выполняет обновление состояния. После \(+1\),
\begin{equation}
(\alpha_u,\alpha_d)\longrightarrow(1,0),
\end{equation}
тогда как после \(-1\),
\begin{equation}
(\alpha_u,\alpha_d)\longrightarrow(0,1).
\end{equation}
Затем он возобновляет унитарную эволюцию начиная с состояния после измерения. Если перед нажатием \(M\) пользователь поворачивает прибор, компьютер меняет базис и вычисляет вероятности, соответствующие новой оси.

Само аппаратное обеспечение не требует ничего специфически квантового. Детерминированного псевдослучайного алгоритма достаточно, чтобы породить последовательность, выглядящую для пользователя случайной. Для одного спина эта классическая имитация может быть операционально неотличима от лабораторного эксперимента.

\begin{classroomqa}
\audiencequestion{Какова роль кнопки измерения? Имеет ли значение прошедшее до её нажатия время?}
\lecturerresponse{Кнопка отмечает момент, в который симулятор производит выбор по распределению Борна и обновляет сохранённые амплитуды. Пока идёт ожидание, амплитуды эволюционируют под действием гамильтониана. Если гамильтониан равен нулю, со временем ничего не меняется; при наличии магнитного поля смоделированный спин прецессирует до измерения.}
\end{classroomqa}

\begin{classroomqa}
\audiencequestion{Кто решает, когда нажать кнопку и как ориентировать прибор?}
\lecturerresponse{В этой модели эти настройки задаёт экспериментатор. Симулятор непрерывно отслеживает состояние, а затем, когда нажимается кнопка, использует выбранную ось и момент измерения. Позднее мы вернёмся к предположению, что настройки можно выбирать независимо от подготовленной пары.}
\end{classroomqa}

\section{Два спина: где локальная симуляция терпит неудачу}

С двумя спинами в состоянии-произведении всё по-прежнему просто. Дайте Алисе один компьютер, хранящий амплитуды \(\alpha_u,\alpha_d\), а Бобу — другой, хранящий \(\beta_u,\beta_d\). Каждый компьютер эволюционирует и измеряет свой спин. Связь не нужна, потому что совместные коэффициенты разлагаются в произведение:
\begin{equation}
\Psi(a,b)=\alpha_a\beta_b.
\end{equation}

Произвольное состояние двух спинов требует четырёх коэффициентов,
\begin{equation}
\alpha_{uu},\qquad
\alpha_{ud},\qquad
\alpha_{du},\qquad
\alpha_{dd},
\qquad
\sum_{a,b}|\alpha_{ab}|^2=1.
\label{eq:l8-four-amplitudes}
\end{equation}
Один центральный компьютер может хранить эти числа и эволюционировать их по четырёхмерному уравнению Schr\"odinger. Когда любой из наблюдателей проводит измерение, центральная машина выбирает исход согласно правильной маргинальной вероятности, отбрасывает несовместимые компоненты, перенормирует оставшиеся компоненты и продолжает эволюцию из условного состояния.

\begin{figure}[htbp]
\centering
\includegraphics[width=0.88\textwidth]{lecture_08_figure_05.png}
\caption{Центральный симулятор хранит четыре амплитуды произвольного состояния двух спинов и остаётся соединённым с обеими станциями измерения.}
\lectureframe{00:37:47}
\end{figure}

Эта конструкция безупречно имитирует запутанность, пока центральный процессор остаётся соединённым с обеими сторонами. Трудности начинаются, когда четыре амплитуды копируют в две машины, разносят их, перерезают провода и позволяют каждой машине использовать собственный генератор случайных чисел. Для синглета,
\begin{equation}
\ket{\Psi^-}
=
\frac{1}{\sqrt2}
\left(
\ket{u}_A\ket{d}_B-\ket{d}_A\ket{u}_B
\right),
\label{eq:l8-singlet}
\end{equation}
измерения вдоль одинаковых осей должны давать противоположные результаты. Независимые случайные выборы могут вместо этого дать \(++\) или \(--\), так что корреляция теряется, хотя обе машины изначально имели копию совместных коэффициентов.

\begin{classroomqa}
\audiencequestion{Могут ли две машины пользоваться одним общим источником случайных чисел до того, как их разнесут?}
\lecturerresponse{Общая случайность может обеспечить заданную корреляцию для одной фиксированной оси измерения. Более глубокая проблема Белла требует одной локальной схемы, которая воспроизводила бы квантовые корреляции для каждой независимо выбранной пары осей. Заранее загруженные общие переменные входят в число белл-локальных моделей, и ни одна такая модель не воспроизводит полную квантовую корреляционную функцию.}
\end{classroomqa}

\begin{classroomqa}
\audiencequestion{Являются ли четыре амплитуды независимыми?}
\lecturerresponse{Это четыре координаты одного нормированного совместного состояния, а не четыре независимых локальных состояния. В запутанном состоянии их нельзя записать как \(\alpha_a\beta_b\). После того как выбран локальный результат, оставшиеся совместные коэффициенты должны быть когерентно обновлены; две разъединённые копии не выполняют автоматически согласованных обновлений.}
\end{classroomqa}

\subsection{Что означают скрытые провода}

Центральный генератор случайных чисел, соединённый с обеими станциями, мог бы координировать результаты. В обычном классическом симуляторе эти провода были бы также каналами связи. Боб мог бы намеренно послать Алисе бит, и если бы связь была мгновенной, машина нарушала бы локальность.

Квантовая система не предоставляет Бобу такой свободы. Допустимые для него локальные преобразования унитарны, и доказательство с редуцированной матрицей плотности уже показало, что ни одно такое преобразование не меняет маргинальную статистику Алисы. Поэтому провода симулятора — признак того, что обычная классическая схема учёта имеет неверную структуру, а не свидетельство того, что в природе есть доступные сверхсветовые кабели.

\begin{classroomqa}
\audiencequestion{Представляют ли провода, необходимые классическому симулятору, физическую нелокальность?}
\lecturerresponse{Они представляют ту дополнительную координацию, которая нужна наивной классической модели. Если бы это были обычные пригодные к использованию провода, они допускали бы передачу сигналов и не моделировали бы квантовый эксперимент. Более надёжный вывод состоит в том, что локальная классическая симуляция и квантовая механика имеют разные логические структуры, а не в том, что обнаружены скрытые провода связи.}
\end{classroomqa}

\begin{classroomqa}
\audiencequestion{Если Алиса и Боб получают по случайному биту, но биты идеально коррелированы, что в этом странного?}
\lecturerresponse{В одной заранее устроенной корреляции нет ничего странного; её могут создать одинаковые конверты или общие случайные затравки. Неклассическое содержание проявляется, когда Алиса и Боб независимо выбирают между несовместимыми настройками, а возникающие корреляции нарушают неравенство Белла, тогда как каждый локальный результат остаётся случайным.}
\end{classroomqa}

В этом мысленном эксперименте есть и предпосылка, которую стоит ясно обозначить: настройки измерения рассматриваются как независимые выборы. Алиса и Боб решают, когда нажимать свои кнопки и как ориентировать свои приборы, без того чтобы эти выборы были тайно предопределены переменными источника. Как только наблюдатели и механизмы их выбора включаются внутрь одной замкнутой физической системы, философский язык свободного выбора становится деликатным. Операциональные тесты Белла формулируют нужное условие как независимость настроек измерения и организуют генерируемые извне настройки так, чтобы они были пространственноподобно разделены. Лекция оставляет метафизику открытой, сохраняя экспериментальную предпосылку.

\begin{classroomqa}
\audiencequestion{Создаёт ли включение Алисы и Боба внутрь физической системы лазейку?}
\lecturerresponse{Это выявляет одно предположение: настройки не должны быть заранее скоррелированы со скрытыми переменными, подготавливающими пару. О свободе воли можно спорить, но проверяемое утверждение — это независимость настроек измерения. Отказаться от неё логически возможно, хотя это смещает объяснительную нагрузку, а не даёт обычную локальную симуляцию.}
\end{classroomqa}

\section{От спинов к частице на прямой}

На примере спинов наше обсуждение прошло от конечномерных векторов через запутанность и локальность. Теперь расширим пространство Гильберта. Рассмотрим одну частицу, движущуюся по бесконечной прямой. Собственное состояние положения обозначается
\begin{equation}
\ket{x},
\end{equation}
и означает, что частица имеет определённое положение \(x\). Эти кеты образуют непрерывно параметризованный базис. Общее состояние является их суперпозицией, а его волновая функция в координатном представлении имеет вид
\begin{equation}
\psi(x)=\braket{x}{\Psi}.
\label{eq:l8-position-wavefunction}
\end{equation}

\begin{figure}[htbp]
\centering
\includegraphics[width=0.88\textwidth]{lecture_08_figure_06.png}
\caption{Волновая функция в координатном представлении как проекция абстрактного состояния на базисный кет \(\ket{x}\).}
\lectureframe{00:56:51}
\end{figure}

Положение используется как метка полного базиса. Это не означает ни что классическое состояние можно задать одним лишь положением, ни что частица одновременно имеет определённые положение и импульс. Это означает, что любой вектор квантового состояния можно разложить по координатному базису. Импульс будет появляться как оператор, действующий на эти коэффициенты в координатном представлении.

Правило Борна теперь даёт \emph{плотность} вероятности,
\begin{equation}
\rho(x)=|\psi(x)|^2.
\label{eq:l8-probability-density}
\end{equation}
Для непрерывной переменной вероятность одной точно заданной точки обычно равна нулю. Вероятность обнаружить частицу в интервале равна
\begin{equation}
P(x_1<x<x_2)
=
\int_{x_1}^{x_2}|\psi(x)|^2\,dx,
\end{equation}
а условие нормировки требует
\begin{equation}
\int_{-\infty}^{\infty}|\psi(x)|^2\,dx=1.
\label{eq:l8-continuous-normalization}
\end{equation}

\begin{classroomqa}
\audiencequestion{Как одно только положение может образовывать полный базис, если классической частице нужны и положение, и импульс?}
\lecturerresponse{Полнота здесь относится к представлению квантового вектора. Коэффициенты \(\psi(x)\) для всех \(x\) задают состояние, а информация об импульсе закодирована в их фазах и пространственном изменении. Положение и импульс — разные базисы, позже связанные преобразованием Фурье; это не одновременные классические метки.}
\end{classroomqa}

\begin{classroomqa}
\audiencequestion{Является ли \(|\psi(x)|^2\) вероятностью обнаружить частицу в точной точке \(x\)?}
\lecturerresponse{Это плотность. Умножение на малый интервал \(dx\) даёт приближённую вероятность для этого интервала, а интегрирование — точную вероятность для конечной области. Плотность может быть отлична от нуля, хотя вероятность, приписанная одной изолированной точке, равна нулю.}
\end{classroomqa}

\section{Непрерывная ортонормированность и дельта-функция}

Для дискретного ортонормированного базиса,
\begin{equation}
\braket{a}{a'}=\delta_{aa'}.
\end{equation}
Для непрерывного базиса аналог имеет вид
\begin{equation}
\braket{x}{x'}=\delta(x-x'),
\label{eq:l8-delta-normalization}
\end{equation}
где дельта-функция Дирака обращается в нуль при \(x\neq x'\) и имеет единичную площадь:
\begin{equation}
\int_{-\infty}^{\infty}\delta(x-x')\,dx'=1.
\end{equation}
Её можно наглядно представить как узкий пик ширины \(\epsilon\), высоты \(1/\epsilon\) и фиксированной площади, с последующим переходом к пределу \(\epsilon\to0\).

\begin{figure}[htbp]
\centering
\includegraphics[width=0.88\textwidth]{lecture_08_figure_07.png}
\caption{Изображение дельта-функции Дирака на доске в виде узкого пика: ширина \(\epsilon\), высота \(1/\epsilon\) и единичная площадь.}
\lectureframe{01:01:57}
\end{figure}

Эта запись — сокращённое обозначение в терминах распределений. Собственный кет положения не является нормируемым вектором в обычном гильбертовом пространстве квадратно-интегрируемых функций; это обобщённый вектор, чья дельта-нормировка имеет смысл под знаком интеграла. Основное рабочее правило имеет вид
\begin{equation}
\int_{-\infty}^{\infty}
\delta(x-x')f(x')\,dx'
=f(x).
\label{eq:l8-delta-sampling}
\end{equation}
Дельта выделяет значение, при котором её аргумент обращается в нуль.

Теперь запишем непрерывное разложение
\begin{equation}
\ket{\Psi}
=
\int_{-\infty}^{\infty}
\psi(x')\ket{x'}\,dx'.
\label{eq:l8-continuous-expansion}
\end{equation}
Действуя \(\bra{x}\), получаем
\begin{align}
\braket{x}{\Psi}
&=
\int dx'\,\psi(x')\braket{x}{x'} \notag\\
&=
\int dx'\,\psi(x')\delta(x-x') \notag\\
&=
\psi(x).
\label{eq:l8-expansion-consistency}
\end{align}
Следовательно, коэффициент в разложении и проекция \(\braket{x}{\Psi}\) — одна и та же функция, в точности как в дискретном базисе.

\begin{figure}[htbp]
\centering
\includegraphics[width=0.88\textwidth]{lecture_08_figure_08.png}
\caption{Непрерывное разложение состояния и проверка его согласованности с помощью \(\braket{x}{x'}=\delta(x-x')\).}
\lectureframe{01:07:00}
\end{figure}

\begin{classroomqa}
\audiencequestion{Играет ли \(\psi(x')\) ту же роль, что и амплитуды в конечномерном векторе?}
\lecturerresponse{Да. Метка базиса \(a'\) становится непрерывной меткой \(x'\), сумма становится интегралом, а коэффициент \(\psi(a')\) превращается в функцию \(\psi(x')\). Лежащее в основе утверждение о векторном пространстве не меняется.}
\end{classroomqa}

\begin{classroomqa}
\audiencequestion{Является ли дельта-функция Дирака в самом деле функцией с бесконечным значением в одной точке?}
\lecturerresponse{Этот образ полезен, но не буквальный. Строго говоря, дельта — это распределение, определяемое своим действием под знаком интеграла, уравнение~\eqref{eq:l8-delta-sampling}. Узкий пик единичной площади — это последовательность обычных функций, которая приводит к тому же правилу выборки.}
\end{classroomqa}

\section{Функции как векторы}

Чтобы увидеть скалярное произведение, сперва напомним дискретное вычисление:
\begin{align}
\braket{\psi}{\phi}
&=
\sum_{a',a}
\psi^*(a')\braket{a'}{a}\phi(a) \notag\\
&=
\sum_a\psi^*(a)\phi(a).
\label{eq:l8-discrete-inner-product}
\end{align}
Комплексно-сопряжённые коэффициенты принадлежат бра, а ортонормированность сводит двойную сумму к одинарной.

\begin{figure}[htbp]
\centering
\includegraphics[width=0.88\textwidth]{lecture_08_figure_09.png}
\caption{Дискретное скалярное произведение, разложенное по базису до замены сумм интегралами.}
\lectureframe{01:11:11}
\end{figure}

Заменим метки базиса положениями, а сумму — интегралом:
\begin{equation}
\braket{\psi}{\phi}
=
\int_{-\infty}^{\infty}\psi^*(x)\phi(x)\,dx.
\label{eq:l8-continuous-inner-product}
\end{equation}
При перестановке векторов результат комплексно сопрягается,
\begin{equation}
\braket{\phi}{\psi}
=
\braket{\psi}{\phi}^*,
\end{equation}
а норма равна
\begin{equation}
\|\psi\|^2=\braket{\psi}{\psi}
=\int |\psi(x)|^2\,dx.
\end{equation}

Следовательно, функции являются векторами: их можно складывать, умножать на комплексные скаляры, нормировать и соединять скалярным произведением. Физические волновые функции на бесконечной прямой квадратно-интегрируемы, поэтому функции, неограниченно растущие на бесконечности, исключаются. Обобщённые базисные функции \(\delta(x-x_0)\) лежат чуть вне этого гильбертова пространства, но остаются незаменимыми как базисные векторы в смысле распределений.

\begin{classroomqa}
\audiencequestion{Должны ли непомеченные пределы интегрирования быть \(-\infty\) и \(+\infty\)?}
\lecturerresponse{Да, для бесконечной прямой. Эти пределы также объясняют, почему нормируемые волновые функции должны достаточно быстро убывать: иначе интеграл полной вероятности не существовал бы.}
\end{classroomqa}

\begin{classroomqa}
\audiencequestion{В одном из выражений на доске отсутствовала вертикальная черта. Меняет ли это смысл?}
\lecturerresponse{Это была лишь описка в обозначениях. Перекрытие базисных векторов имеет вид \(\braket{a'}{a}\), с бра слева и кетом справа. Восстановление черты проясняет, что это скалярное произведение, а не неоднозначная строка или столбец.}
\end{classroomqa}

\begin{classroomqa}
\audiencequestion{Почти всё большую часть времени запутано с окружающей средой?}
\lecturerresponse{В реалистичных макроскопических условиях — да. Системы непрерывно бомбардируются фотонами, молекулами воздуха и другими степенями свободы, так что фазовая информация рассеивается в окружающую среду. Эта запутанность с окружающей средой играет центральную роль в декогеренции и в том, почему интерференцию между макроскопическими альтернативами трудно наблюдать.}
\end{classroomqa}

\begin{classroomqa}
\audiencequestion{Изменила ли математическая теория распределений физику или лишь сделала дельта-функцию респектабельной?}
\lecturerresponse{Её непосредственная заслуга состояла в том, чтобы сделать строгими преобразования, которыми физики уже пользовались. Она дала точный язык для обобщённых функций и уменьшила споры с математиками, хотя сама по себе не разрешила каждую трудность, связанную с распределениями в квантовой теории поля.}
\end{classroomqa}

\section{Линейные операторы, действующие на волновые функции}

Линейный оператор переводит одну функцию в другую, соблюдая суперпозицию:
\begin{equation}
L(c_1\psi+c_2\phi)
=
c_1L\psi+c_2L\phi.
\end{equation}
Наблюдаемая величина должна представляться эрмитовым оператором. В используемых здесь обозначениях условие имеет вид
\begin{equation}
\braket{\phi}{L\psi}
=
\braket{L\phi}{\psi},
\label{eq:l8-hermitian-condition}
\end{equation}
или, что эквивалентно,
\begin{equation}
\braket{\phi}{L\psi}
=
\braket{\psi}{L\phi}^*.
\end{equation}
Для неограниченных операторов это утверждение включает область определения и граничные условия; в вычислениях ниже предполагаются такие волновые функции, для которых граничные члены исчезают.

\subsection{Положение как умножение}

Простейший оператор умножает волновую функцию на координату:
\begin{equation}
(x\psi)(x)=x\psi(x).
\label{eq:l8-position-operator}
\end{equation}
Он линеен, и
\begin{align}
\braket{\phi}{x\psi}
&=
\int dx\,\phi^*(x)x\psi(x), \notag\\
\braket{x\phi}{\psi}
&=
\int dx\,[x\phi(x)]^*\psi(x) \notag\\
&=
\int dx\,\phi^*(x)x\psi(x),
\end{align}
поскольку \(x\) вещественно. Следовательно, умножение на \(x\) эрмитово и представляет наблюдаемую положения.

\begin{figure}[htbp]
\centering
\includegraphics[width=0.88\textwidth]{lecture_08_figure_10.png}
\caption{Проверка эрмитовости оператора положения, рассматриваемого как умножение волновой функции на вещественную переменную \(x\).}
\lectureframe{01:26:47}
\end{figure}

\begin{classroomqa}
\audiencequestion{В \(x\psi(x)\), является ли \(x\) оператором или просто вещественным числом?}
\lecturerresponse{В каждой точке это вещественное значение координаты, но оператором является правило, которое переводит всю функцию \(\psi(x)\) в новую функцию \(x\psi(x)\). Именно то, что эта величина вещественна, делает два скалярных произведения в проверке эрмитовости одинаковыми.}
\end{classroomqa}

\subsection{Дифференцирование антиэрмитово}

Рассмотрим производную
\begin{equation}
D=\frac{d}{dx}.
\end{equation}
Она линейна, но интегрирование по частям даёт
\begin{align}
\braket{\phi}{D\psi}
&=
\int dx\,\phi^*(x)\frac{d\psi}{dx} \notag\\
&=
\left[\phi^*(x)\psi(x)\right]_{-\infty}^{\infty}
-
\int dx\,\frac{d\phi^*}{dx}\psi(x) \notag\\
&=
-\braket{D\phi}{\psi},
\label{eq:l8-derivative-adjoint}
\end{align}
где граничный член исчезает. Следовательно,
\begin{equation}
D^\dagger=-D.
\end{equation}
Производная антиэрмитова, а не эрмитова.

\begin{figure}[htbp]
\centering
\includegraphics[width=0.88\textwidth]{lecture_08_figure_11.png}
\caption{Интегрирование по частям переносит производную между двумя волновыми функциями и вносит решающий минус.}
\lectureframe{01:32:30}
\end{figure}

Умножим на \(-i\):
\begin{equation}
p=-i\frac{d}{dx}.
\label{eq:l8-momentum-natural-units}
\end{equation}
Комплексное сопряжение превращает \(-i\) в \(+i\), тогда как интегрирование по частям даёт ещё один минус. Знаки сокращаются, так что \(p^\dagger=p\). Восстанавливая единицы,
\begin{equation}
p=-i\hbar\frac{d}{dx}.
\label{eq:l8-momentum-operator}
\end{equation}

Среднее значение явно показывает его вещественность:
\begin{align}
\expect{p}_\psi
&=
-i\int dx\,\psi^*(x)\frac{d\psi}{dx} \notag\\
&=
i\int dx\,\frac{d\psi^*}{dx}\psi(x) \notag\\
&=
\expect{p}_\psi^*.
\label{eq:l8-momentum-expectation-real}
\end{align}
Любое число, равное своему комплексно-сопряжённому, вещественно, как и должно быть для среднего значения наблюдаемой величины.

\begin{figure}[htbp]
\centering
\includegraphics[width=0.88\textwidth]{lecture_08_figure_12.png}
\caption{Среднее значение импульса до и после интегрирования по частям; эти две формы являются комплексно-сопряжёнными и равны.}
\lectureframe{01:39:44}
\end{figure}

\begin{classroomqa}
\audiencequestion{Почему оператор, содержащий \(i\), может представлять вещественную наблюдаемую?}
\lecturerresponse{Эрмитовость относится к полному оператору в скалярном произведении, а не к тому, вещественен ли каждый символ в его координатной формуле. Множитель \(i\) компенсирует минус, возникающий при переносе производной интегрированием по частям, делая все средние значения вещественными.}
\end{classroomqa}

\section{Канонический коммутатор}

Ещё одно вычисление показывает, почему оператор в уравнении~\eqref{eq:l8-momentum-operator} заслуживает внимания. В единицах, где \(\hbar=1\), сначала подействуем \(x\), а затем \(p\):
\begin{align}
(px\psi)(x)
&=
-i\frac{d}{dx}\bigl(x\psi(x)\bigr) \notag\\
&=
-i\psi(x)-ix\frac{d\psi}{dx}.
\label{eq:l8-px}
\end{align}
Теперь поменяем порядок:
\begin{equation}
(xp\psi)(x)
=
-ix\frac{d\psi}{dx}.
\label{eq:l8-xp}
\end{equation}
Вычитая уравнение~\eqref{eq:l8-px} из уравнения~\eqref{eq:l8-xp},
\begin{equation}
\bigl([x,p]\psi\bigr)(x)
=
\bigl(xp-px\bigr)\psi(x)
=
i\psi(x).
\end{equation}
Поскольку это верно для любой волновой функции в области определения оператора,
\begin{equation}
[x,p]=i,
\end{equation}
или, после восстановления единиц,
\begin{equation}
[x,p]=i\hbar.
\label{eq:l8-canonical-commutator}
\end{equation}

\begin{figure}[htbp]
\centering
\includegraphics[width=0.88\textwidth]{lecture_08_figure_13.png}
\caption{Завершённое операторное рассуждение: если \(([x,p]-i)\psi=0\) для любой волновой функции, то \([x,p]=i\).}
\lectureframe{01:46:24}
\end{figure}

Классическим аналогом является
\begin{equation}
\{x,p\}_{\mathrm{PB}}=1.
\end{equation}
Историческое рассуждение Дирака шло в направлении, противоположном настоящему выводу. Он исходил из алгебраического сходства между антисимметрическими скобками Пуассона и коммутаторами, постулировал соответствие
\begin{equation}
[A,B]\longleftrightarrow i\hbar\{A,B\}_{\mathrm{PB}},
\end{equation}
и заключил, что оператор импульса должен удовлетворять каноническому коммутатору с координатой. Здесь мы получили \(-i\,d/dx\), устранив неэрмитовость производной, и лишь затем обнаружили ту же самую алгебру.

\begin{classroomqa}
\audiencequestion{Доказывает ли один лишь коммутатор, что \(-i\hbar\,d/dx\) — это физический импульс?}
\lecturerresponse{Это веское свидетельство, но ещё не полное физическое отождествление. Нам всё ещё нужно показать, что оператор порождает пространственные трансляции, что его среднее значение сохраняется при наличии трансляционной симметрии и что волновой пакет движется с групповой скоростью \(\expect{p}/m\) в соответствующей нерелятивистской системе.}
\end{classroomqa}

\section{Что дальше}

Переход от спинов к частицам сохранил архитектуру, изменив лишь её масштаб. Базис, индексируемый непрерывным параметром, заменяет конечный список базисных векторов. Суммы превращаются в интегралы, дельты Кронекера — в дельты Дирака, а векторы-столбцы — в квадратично интегрируемые функции. Оператор координаты действует умножением, тогда как оператор импульса — дифференцированием.

Следующие задачи носят физический характер, а не сводятся лишь к обозначениям. Для вероятностей импульса нужны базис импульса и преобразование Фурье. Временная эволюция требует уравнения Schr\"odinger для \(\psi(x,t)\). Затем волновые пакеты свяжут абстрактный оператор \(p=-i\hbar\,d/dx\) с движением: их групповая скорость будет связана с импульсом, делённым на массу, а трансляционная симметрия объяснит сохранение импульса. Поэтому коммутатор \([x,p]=i\hbar\) не является изолированным приёмом. Он служит связующим звеном между уже развитой механикой векторного пространства и квантовой динамикой частиц.
